\setcounter{ex}{0} % Lương Thế Vinh - Đà Nẵng
\subsection{Trắc nghiệm (2 điểm)}
\Opensolutionfile{ans}[ans/ans9] 

%%Câu 1:
\begin{ex}
Điều kiện để biểu thức $\sqrt{x-1}+\sqrt{2-x}$ có nghĩa là
\choice
{\True $1 \le x \le 2$}
{$x \ge 2$}
{$x \ge 1$}
{$x \le 2$}
\end{ex}

%%Câu 2:
\begin{ex}
Cho hai đường tròn $(O_1)$ và $(O_2)$ lần lượt có bán kính bằng $3$ và $5$, đồng thời $O_1O_2=6$. Khẳng định nào dưới đây là đúng?
\choice
{$(O_1)$ và $(O_2)$ không cắt nhau}
{$(O_1)$ và $(O_2)$ tiếp xúc trong}
{\True $(O_1)$ và $(O_2)$ có hai điểm chung}
{$(O_1)$ và $(O_2)$ tiếp xúc ngoài}
\end{ex}

%%Câu 3:
\begin{ex}
Hệ phương trình $\heva{&2(x-1)-3y=3\\
&3x+2(y-1)=-1}$ có nghiệm là $(x;y)=(a;b)$. Giá trị của biểu thức $T=a+b$ bằng
\choice
{$2$}
{\True $0$}
{$-2$}
{$3$}
\end{ex}

%%Câu 4:
\begin{ex}
Cho tam giác $ABC$ có $AB=3, BC=4, AC=5$. Độ dài đường cao $AH$ bằng bao nhiêu?
\choice
{\True $\dfrac{12}{5}$}
{$\dfrac{24}{5}$}
{$12$}
{$24$}
\end{ex}

%%Câu 5:
\begin{ex}
Trong các phương trình dưới đây, đâu là phương trình bậc nhất hai ẩn?
\choice
{$2x+y=3-z$}
{$3x=2y^2$}
{$x-2y=3x(1-y)$}
{\True $3x+y=2$}
\end{ex}

%%Câu 6:
\begin{ex}
Cho hai đường tròn $(O_1)$ và $(O_2)$ cùng có bán kính bằng $2024$. Biết rằng đường tròn này đi qua tâm của đường tròn kia và $A$ nằm trên $(O_1)$ sao cho $A,O_1,O_2$ thẳng hàng theo thứ tự đó. Tính độ dài tiếp tuyến kẻ từ $A$ đến $(O_2)$ (kết quả làm tròn hai chữ số thập phân).
\choice
{\True $3\,505{,}67$}
{$2\,862{,}37$}
{$4\,498{,}89$}
{$6\,362{,}39$}
\end{ex}

%%Câu 7:
\begin{ex}
Cho $a$ và $b$ là hai số thực dương. Nếu $a>b$ thì trong các khẳng định dưới đây, có bao nhiêu khẳng định đúng?
	\begin{listEX}[4]
		\item[i.] $2a+1>2b+1$.
		\item[ii.] $-3a-2>-3b-2$.
		\item[iii.] $\dfrac{2}{a}<\dfrac{2}{b}$.
		\item[iv.] $a^2\geq b^2$.
	\end{listEX}
\choice
{$4$}
{\True $3$}
{$2$}
{$1$}
\end{ex}

%%Câu 8:
\begin{ex}
Nghiệm nguyên lớn nhất của bất phương trình $x(2-x)+2x^2<8-x+x(x-1)$ bằng
\choice
{$2$}
{\True $1$}
{$0$}
{Không tồn tại}
\end{ex}

%Câu 9
\begin{ex}
Tìm góc nhọn $\alpha$, biết $\sin \alpha=\cos \alpha$
\choice
{$60^\circ$}
{\True $45^\circ$}
{$30^\circ$}
{$15^\circ$}
\end{ex}

%Câu 10 
\begin{ex}
Một tòa nhà cao tầng có bóng trên mặt đất dài 272 m, cùng thời điểm đó một cột đèn cao 7 m có bóng trên mặt đất dài 14 m. Tòa nhà đó có bao nhiêu tầng, biết rằng mỗi tầng cao 3,4 m?
\choice
{$42$ tầng}
{$38$ tầng}
{$35$ tầng}
{\True $40$ tầng}
\end{ex}


\Closesolutionfile{ans}
\Opensolutionfile{ansbook}[ans/ansbook9] 
\setcounter{bt}{0} 
\subsection{Tự luận (8 điểm)}
%Bài 1
\begin{bt}[1,0 điểm] $ $
\begin{enumerate}
    \item Tính: $\sqrt{25}\cdot \sqrt{16}+20:\sqrt[3]{-8}$
    \item Rút gọn biểu thức: $A=\left(\dfrac{\sqrt{14}-\sqrt{7}}{1-\sqrt{2}}+\dfrac{\sqrt{15}-\sqrt{5}}{1-\sqrt{3}}\right): \dfrac{1}{\sqrt{7}-\sqrt{5}}$
\end{enumerate}
\loigiai{
a) $10$ \hspace{2cm} b) $A=-2$
}
\end{bt}

%Bài 2
\begin{bt}[2,0 điểm] $ $
\begin{enumerate}
    \item Cho biểu thức: $\mathrm{Q}=\dfrac{2}{2+\sqrt{x}}+\dfrac{1}{2-\sqrt{x}}+\dfrac{2 \sqrt{x}}{x-4}$
    \begin{enumerate}
        \item[i)] Rút gọn $\mathrm{Q}$
        \item[ii)] Tìm các giá trị nguyên $x$ để $\mathrm{Q}$ nhận giá trị nguyên.
    \end{enumerate}
    \item Giải phương trình: $\dfrac{x-1}{x-5}-\dfrac{4 x}{x+5}=\dfrac{9 x-5}{x^2-25}$
\end{enumerate}
\loigiai{
a) i) Q$=\dfrac{3}{\sqrt{x}+2}$, ii) $x=1$ \hfill b) $x=0$ \hfill
}
\end{bt}

%Bài 3
\begin{bt}[1,0 điểm] $ $
Một người đứng cách một cây xanh một khoảng 30 m. Mắt người đó nhìn lên đỉnh ngọn cây và tạo với phương nằm ngang một góc $20^{\circ}$. Mắt người cách mặt đất 1,6 m. Tính chiểu cao của cây. (Tính chính xác đến dm).
\begin{center}
\begin{tikzpicture}[scale=0.7, font=\small, >=stealth]
\coordinate (N) at (0,0);
\coordinate (M) at (0,1.5); 
\coordinate (A) at (15,0); 
\coordinate (H) at (15,1.5);
\coordinate (B) at (15,7); 
\node[anchor=south west] at (11.5,-0.2) {\includegraphics[height=5cm]{hinh/De9-tree.png}};
\draw[thick] (N) -- (A);
\node[below] at ($(N)!0.5!(A)$) {30 m};
\draw[dashed, thick] (N) -- (M);
\node[right] at ($(N)!0.5!(M)$) {1,6 m};
\draw[dashed, thick] (M) -- (B);
\draw[dashed] (M) -- (H);
\draw[thick] (A) -- (B);
\draw[thick] (2,1.5) arc[start angle=0, end angle=20, radius=2];
\node at (2.5,1.9) {$20^\circ$};
\foreach \p/\pos in {A/below,B/above right,H/right,M/left,N/below}
      \filldraw[fill=black] (\p) circle(0.5pt) node[\pos] {$\p$};
\node[anchor=south west] at (-0.4,-0.15) {\includegraphics[height=1cm]{hinh/De9-man.png}};
\end{tikzpicture}
\end{center}
\loigiai{Vận tốc dự định $40 \text{ km/h}$ ; quãng đường AB dài $120 \text{ km}$.}
\end{bt}

%Bài 4
\begin{bt}[3,0 điểm]
Từ điểm $M$ nằm ngoài đường tròn $(O, R)$ kẻ hai tiếp tuyến $MB,\, MC$ ($B, C$ là hai tiếp điểm). $OM$ cắt $BC$ tại $H$.
\begin{listEX}
    \item Giả sử $OM=2R$. Tính $BC$ theo $R$.
    \item Kẻ đường kính $BD$ của $(O)$, $MD$ cắt $(O)$ tại $K$. Chứng minh rằng $MK\cdot MD=MH\cdot MO$.
    \item Chứng minh rằng $HC$ là phân giác của góc $KHD$.
\end{listEX}
\loigiai{\immini{
	\begin{listEX}
		\item $BC=R\sqrt{3}$.
		\item Chứng minh $\triangle MCK \backsim \triangle MDC$ từ đó suy ra $MC^2=MK. MD$.
		
		Mặt khác $MC^2=MO\cdot MH$. Từ đó suy ra điều phải chứng minh.
		\item Chứng minh $\widehat{OHD}=\widehat{ODM}=\widehat{KHM}$ thông qua các tam giác đồng dạng.
		\end{listEX}}
	{
	\begin{tikzpicture}[line width=0.8pt]
		\def\R{2}
		\def\A{6}
		\path 
		(0,0) coordinate (O)
		(\A,0) coordinate (M)
		($(O)!0.5!(M)$) coordinate (I)
		;
		\path[name path=ciro] (O) circle (\R);
		\path[name path=cirom] (I) let \p1=($(I)-(O)$) in circle ({veclen(\x1,\y1)});
		\path[name intersections={of=ciro and cirom, by={B,C}}];
		\coordinate (H) at (intersection of O--M and B--C);
		\coordinate (D) at ($(B)!2!(O)$);
		\path[name path=md] (D)--(M);
		\path[name intersections={of=ciro and md,by={D,K}}];
		\draw (B)--(C) (M)--(D);
		\draw
		(O) circle (\R)
		(B)--(O)--(C) (O)--(M)
		(B)--(M)
		(M)--(C) (O)--(D)--(C)--(K)--(H)--(D)
		;
		\foreach \x/\g in{O/-180,M/0,B/60,C/-60,H/45,D/-90,K/90}
		\draw[fill=white] (\x) circle (1.5pt) +(\g:.35) node{$\x$};
	\end{tikzpicture}
	}
}
\end{bt}

%Bài 5
\begin{bt}[1,0 điểm] Một ô tô đi từ A đến B với vận tốc và thời gian dự định. Nếu người đó tăng vận tốc thêm $20 \mathrm{~km} / \mathrm{h}$ thì đến B sớm hơn dự định 1 giờ Nếu người đó giảm vận tốc $10 \mathrm{~km} / \mathrm{h}$ thì đến B muộn hơn 1 giờ. Tính vận tốc dự định và độ dài quãng đường AB 
\loigiai{
Gọi $x$ (km/h) và $y$ (h) lần lượt là vận tốc và thời gian dự định ($x > 10 ; y > 1$).

Từ giá thiết ta có hệ phương trình:
$\begin{cases}
    (x+20)(y-1)=xy \\
    (x-10)(y+1)=xy
\end{cases} $
    \item Giải hệ PT: $x=40$ (nhận) ; $y=3$ (nhận).}
\end{bt}

\Closesolutionfile{ansbook}